\section{GPU reference orbit backend}
\label{subsec:ref-orbit-gpu}

In many numerical algorithms, correctness and stability depend on arithmetic
precision far exceeding that of standard IEEE~754 floating-point formats. The
Mandelbrot reference orbits is a canonical example: small rounding errors
introduced early in the iteration can grow exponentially, eventually corrupting
orbit classification, perturbation terms, or bailout logic. While
arbitrary-precision libraries like MPIR can provide the required accuracy
(\cref{sec:ref-orbit-calc}), their
performance is often insufficient when reference orbits must be computed
repeatedly or at very high precision. This creates a tension between
numerical fidelity and throughput.

To resolve this tension, high-precision arithmetic operations---in particular
multiplication (\cref{sec:ntt-multiply}), addition, and subtraction (\cref{sec:hp-add}) of large significands---must be
implemented with both mathematical rigor and architectural efficiency. Addition
and subtraction are dominated by carry propagation across hundreds or thousands
of limbs, while multiplication scales quadratically unless asymptotically faster
algorithms are employed. For precisions relevant to deep zoom reference orbits,
naive limb-by-limb multiplication quickly becomes the dominant cost,
overwhelming all other parts of the computation.

\begin{figure}[!htbp]
  \centering
  \includegraphics[width=0.9\linewidth]{gpu-reference-orbit.png}
  \caption[View \#30 at zoom $10^{114{,}514}$ via GPU backend]{View \#30 rendered at zoom factor \(10^{114,514}\) using the GPU
  reference orbit backend.  Requires ~73s to render on RTX 5090.}
  \label{fig:gpu-reference-orbit}
\end{figure}

\begin{figure}[!htbp]
  \centering
  \includegraphics[width=0.9\linewidth]{view32.png}
  \caption[View \#32 at zoom $10^{244{,}240}$]{View \#32 rendered at zoom factor \(10^{244,240}\), roughly twice
  the exponent depth of View~\#30.  Reaching this depth exercises both the
  NTT multiplication pipeline and the high-precision addition/subtraction
  kernels at full scale.}
  \label{fig:view32}
\end{figure}

\subsection{Why Fast Multiplication Matters for Reference Orbits}

Generating many fractal reference orbits typically involves iterating a
recurrence relation (e.g.\ $z_{n+1} = f(z_n)$) at very high precision, often for
thousands of iterations per reference point
(\cref{fig:gpu-reference-orbit}).Each iteration requires multiple
high-precision multiplications and additions, and the total cost grows linearly
with iteration count and superlinearly with precision. Even modest speedups in
the core arithmetic therefore compound over the full computation.

By using an NTT-based multiplication strategy, large significand products can be
computed in $O(n \log n)$ time instead of $O(n^2)$, shifting the performance
profile of high-precision arithmetic into a regime where GPUs can be effectively
utilized. When combined with carefully engineered carry propagation and
normalization steps, this enables high-precision multiply/add/subtract
operations that are fast enough to make reference orbit computation practical at
scales that would otherwise be prohibitive (\cref{fig:view32}).

\subsection{System Architecture and Execution Context}
\label{subsec:ref-orbit-architecture}

The GPU reference orbit backend is a stateful, precision-specialized execution
engine that computes authoritative Mandelbrot reference orbits entirely on the
GPU. The backend advances the recurrence
\[
z_{n+1} \leftarrow z_n^2 + c
\]
at fixed high precision by executing multiple iterations per kernel invocation,
while the CPU orchestrates invocation, result collection, and termination.

\paragraph{Division of responsibility.}
The design follows a strict separation of concerns. The CPU is responsible for
selecting a supported precision, initializing GPU resources, invoking the kernel
in bounded iteration batches, and accumulating emitted iteration records into
\code{PerturbationResults}. The GPU owns all numerical state of the orbit,
executes the high-precision arithmetic pipeline, detects escape and periodicity,
and reports authoritative outcomes to the host.

\paragraph{Persistent backend state.}
Reference orbit computation is centered around a persistent \emph{combo} object
(\code{HpSharkReferenceResults<SharkFloatParams>}) that encapsulates all state
required to advance the orbit at a fixed precision. The combo contains:
\begin{itemize}
  \item high-precision orbit-related values (e.g.\ radius threshold and derivative terms),
  \item persistent arithmetic pipeline state for multiplication
        (\code{Multiply}, implementing NTT-based convolution) and
        addition (\code{Add}, implementing parallel-prefix carry
        propagation),
  \item a fixed-size output mailbox (\code{OutputIters},
        \code{Output\-Iter\-Count},
        \code{Max\-Output\-Iters}\,=\,1024),
  \item termination status (\code{PeriodicityStatus}),
  \item and host-only launch and resource plumbing, including a device pointer to
        the GPU-resident state, temporary device buffers, an associated CUDA
        stream, and cached kernel argument arrays.
\end{itemize}
The combo is initialized once, reused across multiple kernel invocations, and
destroyed only after the reference orbit computation completes.

\paragraph{Batch-oriented kernel invocation.}
Each call to \code{InvokeHpSharkReferenceKernel} launches a GPU kernel that
advances the Mandelbrot recurrence by up to a specified number of iterations
(\code{itersToRun}). Within a single invocation, the kernel performs multiple
full iterations internally, combining high-precision squaring and addition on
the GPU. As iterations progress, the kernel emits up to
\code{MaxOutputIters} iteration records into the combo’s output mailbox and
updates \code{OutputIterCount} accordingly.

This batching strategy bounds kernel runtime, enables incremental streaming of
orbit data back to the host, and allows long orbits to be computed without
reinitializing GPU state.

\paragraph{Termination and authority.}
After each invocation, the host inspects the combo’s \code{PeriodicityStatus}
and stops early if the GPU reports escape, detected periodicity, or an error
condition, or if the total requested iteration budget has been reached. In this
model, the GPU backend is authoritative: it owns the orbit state, performs all
high-precision arithmetic exactly with respect to the fixed significand
precision, and determines orbit outcomes without CPU-side recomputation or
correction.

\paragraph{Precision specialization.}
The entire lifecycle of a reference orbit computation is specialized by
\code{SharkFloatParams}. Precision is selected once at startup via
\code{DispatchByPrecision} and remains fixed for the lifetime of the combo
object. This ensures that kernel code, shared-memory usage, and arithmetic
pipelines are fully specialized for the chosen precision, at the cost of
supporting a finite, enumerated set of precisions.

\medskip

This architectural structure defines the execution context for the remainder of
this section. We begin with the underlying \code{HpSharkFloat}
representation and then describe the internal arithmetic pipelines used within
each kernel invocation, including high-precision addition, NTT-based
multiplication, carry propagation, and normalization, explaining how these
components are organized to maintain correctness and performance within the
bounded-chunk execution model.

\subsection{High-Precision Floating-Point Representation}
\label{subsec:hpfloat-model}

The GPU reference orbit backend operates on a custom high-precision
floating-point format implemented by the \code{HpSharkFloat} family of types.
This format is intentionally \emph{not} IEEE~754--compatible. Instead, it is
designed to support exact arithmetic on large significands, deterministic
behavior across CPU and GPU backends, and efficient execution on massively
parallel architectures.

Conceptually, an \code{HpSharkFloat} value represents a real number of the
form
\[
x = s \cdot M \cdot 2^E,
\]
where $s \in \{-1,+1\}$ is the sign, $M$ is a nonnegative integer significand,
and $E$ is a signed integer exponent. The sign, significand, and exponent are
tracked independently, and no implicit normalization or rounding is performed
during arithmetic operations.

\paragraph{Significand representation.}
The significand $M$ is stored as a fixed-size array of 32-bit limbs. The number
of limbs is determined at compile time by the
\code{SharkFloatParams} specialization and remains constant throughout the
lifetime of the computation. As a result, precision is explicit and static:
all arithmetic kernels operate on limb arrays of identical size, and no dynamic
reallocation or precision growth occurs during iteration.

The significand represents a nonnegative integer magnitude. Negative values are
handled by explicit sign logic rather than two’s-complement arithmetic, which
simplifies carry propagation and normalization.

\paragraph{Exponent handling and deferred normalization.}
The exponent $E$ is maintained separately from the significand and updated
explicitly by arithmetic kernels. In contrast to IEEE~754 arithmetic,
normalization is not performed eagerly after every operation. Intermediate
results may remain temporarily denormalized, with exponent adjustments and
carry propagation deferred until well-defined normalization phases.

This separation is essential for GPU execution. Immediate normalization would
introduce serial dependencies and global synchronization at every arithmetic
step, whereas deferred normalization allows large batches of arithmetic to be
performed in parallel before a single, optimized carry-propagation pass.

\paragraph{Normalization and canonical form.}
A value is considered to be in canonical form when:
\begin{itemize}
  \item the significand has been fully carry-propagated,
  \item the most significant limb is nonzero (except for the zero value),
  \item the exponent reflects the exact binary scaling of the significand.
\end{itemize}
Normalization is implemented as a separate phase and is optimized using
parallel-prefix carry propagation techniques, as described later in this
section.

\paragraph{Rounding, exceptional values, and determinism.}
No rounding is performed during addition, subtraction, or multiplication. All
operations are exact with respect to the fixed significand size. If an
intermediate result were to exceed the representable precision, this would
constitute a logic error rather than silently rounded behavior.

The format does not represent NaNs, infinities, or subnormal values. These
concepts are unnecessary for reference orbit computation, where arithmetic
ranges are structurally controlled and exceptional conditions (such as escape)
are handled at the algorithmic level.

Because the precision is fixed, arithmetic is exact, and no data-dependent
rounding occurs, the \code{HpSharkFloat} model provides deterministic,
bitwise-reproducible results across runs and across CPU and GPU backends,
subject only to correct synchronization.

\paragraph{Implications for reference orbit computation.}
This floating-point model is best viewed as high-precision integer arithmetic
with an explicit exponent rather than as an extension of IEEE floating-point.
The design directly enables NTT-based multiplication, parallel carry
propagation, and predictable numerical behavior at extreme precisions, making
it well suited for authoritative reference orbit generation on the GPU.

\subsection{High-Precision Floating-Point Addition on the GPU}
\label{sec:hp-add}

While high-precision multiplication (\cref{sec:ntt-multiply}) dominates the asymptotic cost of many
arbitrary-precision algorithms, addition and subtraction remain critical in
practice. In iterative computations such as reference orbit generation, each
iteration performs multiple additions and subtractions on large significands,
often regardless of whether multiplication is required. At precisions of
hundreds or thousands of limbs, even these ostensibly simple operations become
nontrivial due to long-range carry and borrow propagation.

A naive limbwise implementation resolves carries sequentially: a carry-out from
limb $i$ must be fully determined before limb $i+1$ can be finalized. This
introduces strict serial dependence and $O(N)$ latency, making naive carry
handling incompatible with parallel GPU architectures.
Without reformulation, carry propagation becomes a synchronization
bottleneck that erodes the benefits of fast high-precision multiplication.

In \FractalShark{}'s current reference-orbit pipeline, this add/subtract path
is retained primarily for architectural reasons rather than because isolated
high-precision addition is an overwhelming GPU speedup. Carry propagation
remains an expensive, synchronization-heavy part of the computation. Keeping
the operation on the GPU nevertheless preserves an all-on-device reference-orbit
pipeline: the orbit state stays resident, and a single kernel invocation can
execute multiple full iterations internally instead of returning to the CPU
after each iteration for the add/subtract step and then restarting the kernel.

\subsubsection{Numerical Representation}

We represent a high-precision floating-point number as
\begin{equation}
x = s_x \cdot 2^{e_x} \cdot \sum_{i=0}^{N-1} d_i \, 2^{-w i},
\end{equation}
where $s_x \in \{-1,+1\}$ is the sign, $e_x \in \mathbb{Z}$ is a shared exponent,
$d_i \in \{0,\dots,2^w-1\}$ are fixed-width limbs, $w$ is the limb bit width
(typically $32$ or $64$), and $N$ is the limb count. Arithmetic is performed on
the limb array, while the exponent is tracked separately. The mantissa is thus
treated as a fixed-point integer scaled by an implicit radix $2^{-w}$.

\subsubsection{Problem Statement and Exponent Alignment}

Given two high-precision floating-point numbers
\[
x = s_x \cdot 2^{e_x} \cdot M_x, \qquad
y = s_y \cdot 2^{e_y} \cdot M_y,
\]
the objective is to compute
\[
z = x + y
\]
exactly, or under a well-defined rounding mode, using GPU-parallel execution.

Assuming without loss of generality that $e_x \ge e_y$, we define
\[
\Delta = e_x - e_y
\]
and align the mantissas by shifting
\begin{equation}
M_y' = M_y \cdot 2^{-\Delta}.
\end{equation}
If $\Delta \ge Nw$, the contribution of $y$ is below the representable precision
and the operation degenerates to $z \approx x$. Otherwise, both mantissas are
expressed at the shared exponent $e = e_x$.

\subsubsection{Signed Limbwise Accumulation}

After alignment, addition and subtraction are unified by introducing signed limb
arrays
\[
a_i = s_x \cdot d_i^{(x)}, \qquad
b_i = s_y \cdot d_i^{(y')}.
\]
The raw limbwise sum is then
\begin{equation}
t_i = a_i + b_i,
\end{equation}
which may lie outside the canonical range $[0,2^w)$. This computation is
embarrassingly parallel and maps directly to GPU threads, with each thread
computing one or more limb sums independently.

\subsubsection{Carry Propagation via Function Composition}

Each intermediate sum $t_i = a_i + b_i$ produces a base value $r_i = t_i \bmod 2^w$. However, because $a_i$ and $b_i$ are signed, an incoming carry into limb $i$ may trigger further overflow or underflow, modifying the carry-out. This dependency is captured by a \emph{transfer function} $f_i: S \to S$ over the state space of possible carries $S = \{-1, 0, +1\}$. For each limb $i$, the function $f_i$ is constructed by simulating the carry arithmetic for all three possible input carries: $f_i(c_{\text{in}})$ returns the carry-out $c_{\text{out}}$ that limb $i$ would produce when receiving $c_{\text{in}}$.

The total carry $C_i$ entering limb $i$ is the result of applying the composed transfer functions of all lower-order limbs to the initial carry (zero):
\begin{equation}
C_i = (f_{i-1} \circ f_{i-2} \circ \cdots \circ f_0)(0).
\end{equation}
Function composition is associative --- $(f_c \circ f_b) \circ f_a = f_c \circ (f_b \circ f_a)$ for all $f_a, f_b, f_c$ --- so the transfer functions form a monoid under composition with the identity map as the unit. This allows carry propagation to be computed using a parallel prefix scan, generalizing classical carry-lookahead logic to signed arbitrary-precision arithmetic.

\paragraph{Three-stream batching.}
The full add pipeline requires carry propagation for three independent streams simultaneously: two candidate results for the $A - B + C$ computation (corresponding to the two possible sign-comparison outcomes), and one result for the $D + E$ term. Rather than running three separate scans, the implementation packs all three transfer functions into a single 18-bit struct (\code{PPTransfer3}): three streams $\times$ three input carries $\times$ two bits per output, stored in a \code{uint64\_t}. Composition and prefix operations act on the packed representation, so a single scan pass resolves carries for all three streams at once.

\subsubsection{Single-Pass Parallel Prefix on the GPU}

To evaluate these prefixes efficiently, we employ a single-pass parallel prefix formulation with decoupled look-back \cite{merrill2016single}. The state of each transfer function $f_i$ is encoded compactly in a bit-packed table of the three possible carry inputs, and composition is evaluated directly on those encoded states.

Each thread block computes the local prefix composition of its assigned limb range independently. A lightweight look-back mechanism then resolves inter-block dependencies by querying the aggregate transfer functions of preceding blocks. This design decouples the propagation of the global carry from the strict left-to-right execution order required for prefix consumption. A block can publish its aggregate transfer function immediately upon completion, allowing downstream blocks to accumulate it even if intermediate blocks are still waiting for their own incoming prefixes. In Merrill--Garland terms, a small amount of redundant look-back work is used to hide the latency of global prefix propagation while keeping the scan single-pass and work-efficient.

\subsubsection{Finalization and Normalization}

Once the prefix scan completes, each stream yields the resolved carry $C_i$ at every limb position, and the corrected result limbs are
\begin{equation}
d_i^{(z)} = (r_i + C_i) \bmod 2^w.
\end{equation}
In the $A - B + C$ pipeline, the two candidate streams (\code{extResultTrue} and \code{extResultFalse}) correspond to the two speculative magnitude-ordering hypotheses used to evaluate the three-way sum. Both are fully carry-propagated, and the kernel selects the correct branch by comparing their final carry accumulators; this determines which hypothesis was correct, while the sign itself comes from the precomputed branch state. The third stream ($D + E$) is unconditionally accepted.

A carry beyond the most significant limb increments the exponent,
\[
d_N^{(z)} \neq 0 \;\Rightarrow\; e \leftarrow e + 1,
\]
while leading-zero cancellation triggers renormalization and a corresponding
decrement of the exponent.

\subsubsection{Correctness and Complexity}

The algorithm is mathematically equivalent to exact integer addition of aligned
mantissas,
\[
z = 2^e \cdot (M_x + M_y'),
\]
and the scan operates over associative transfer-function composition, with each
transfer encoding the exact carry effect of its limb. All operations are integer-valued and deterministic, yielding
bitwise-identical results across executions.

For $N$ limbs, the algorithm performs $O(N)$ work with $O(N)$ global memory
traffic. Its practical advantage comes from the single-pass,
communication-avoiding structure: local prefix work is done once per block, and
the extra look-back work is confined to resolving inter-block dependencies
without requiring a separate global scan pass.

\paragraph{Comparison with MPIR/GMP Serial Carry Propagation}

Conventional high-precision libraries such as MPIR and GMP implement addition
using strictly serial carry propagation. After exponent alignment, limbs are
processed sequentially from least to most significant, with each carry-out
immediately applied to the next limb. This approach is optimal for scalar CPUs:
it minimizes memory traffic, exploits tight data locality, and benefits from
branch prediction and instruction-level parallelism. For moderate precisions,
the constant factors are small and the simplicity of the algorithm dominates.

However, the serial carry dependency enforces an $O(N)$ critical path that
limits parallelism. Even when vector instructions are used to
accelerate limbwise addition, the carry chain itself remains sequential and
cannot be overlapped across independent execution units. As a result, MPIR/GMP
addition scales primarily in throughput with clock frequency rather than in
latency with available parallel resources.

In contrast, the GPU-oriented formulation presented here trades a modest
increase in algorithmic complexity for a much more parallel dependency
structure. By expressing carry propagation as an associative prefix operation
and resolving dependencies using a single-pass, decoupled look-back scheme, the
implementation avoids separate global scan passes and allows many partitions to
make progress concurrently. This shift aligns the algorithm with the strengths
of parallel architectures while preserving exact arithmetic
semantics.

The GPU approach does not outperform MPIR/GMP for small limb
counts, where launch overhead and synchronization dominate. Its advantage
emerges at large precisions, where serial carry latency becomes comparable to
or exceeds the cost of multiplication. In this regime, parallel carry
propagation is needed to sustain end-to-end performance in
high-precision iterative workloads.

\paragraph{Role of Merrill--Garland in Parallel Carry Propagation}

The approach adopted here is closely informed by the work of
\cite{merrill2016single} on single-pass parallel prefix scans with decoupled
look-back. Their contribution is not a numerical
algorithm per se, but a general execution strategy for evaluating long,
associative prefix dependencies efficiently on GPUs.

In the context of high-precision addition, carry propagation is naturally
expressed as a prefix composition over per-limb transfer functions. However,
naively mapping this formulation to a GPU using classical parallel scans would
require multiple global synchronization phases or kernel launches, introducing
significant overhead and limiting scalability. The Merrill--Garland framework
addresses this by separating the computation into two logically distinct
components: (1) local prefix evaluation within a thread block, and (2)
resolution of inter-block dependencies through a lightweight, dynamically
resolved look-back mechanism.

The decoupled look-back design allows blocks to complete their local
scan work independently and wait only when an unresolved upstream prefix is
actually needed. This property aligns well with the statistical structure of
carry propagation in large arbitrary-precision additions, where long carry
chains are relatively rare. The algorithm therefore achieves high average
throughput while retaining correctness in worst-case carry scenarios.

It is important to emphasize that Merrill--Garland does not prescribe how carry
information should be represented, nor does it assume a particular numerical
domain. In this work, their scan framework is specialized to the algebra of
carry transfer functions arising from signed limbwise addition. The numerical
semantics---exact integer arithmetic with deterministic behavior---remain
entirely independent of the scan mechanism. In this sense, Merrill--Garland
provides the parallel control structure, while the arithmetic formulation
determines the correctness and stability of the result.

By combining a mathematically exact carry formulation with a single-pass prefix
execution model, the resulting algorithm adapts classical
high-precision arithmetic to modern GPU execution, enabling
carry-dominated operations, traditionally serial, to scale
efficiently on parallel hardware.

\paragraph{Background: parallel prefix scans.}
A \emph{parallel prefix scan} (also called a \emph{prefix sum} or
\emph{scan}) is a parallel primitive: given an array of
$N$~elements and an associative binary operator~$\oplus$, compute the
running ``total'' at every position:
\[
y_i = x_0 \oplus x_1 \oplus \cdots \oplus x_i, \qquad i = 0, \ldots, N-1.
\]
When~$\oplus$ is ordinary addition, this is the classical prefix sum
$(y_0, y_1, \ldots) = (x_0,\; x_0{+}x_1,\; x_0{+}x_1{+}x_2, \ldots)$.
The relevant property is that \emph{any} associative operator works---including
function composition.  For carry propagation, each $x_i$ is a transfer
function $f_i$ mapping an incoming carry to an outgoing carry, and
$\oplus$ is function composition.  The prefix scan therefore computes,
for every limb~$i$, the composed effect of all preceding limbs' carry
behavior---exactly the information needed to resolve the final carry at
each position.

A sequential prefix scan is trivially $O(N)$ work and $O(N)$ depth
(left to right).  A \emph{parallel} prefix scan achieves $O(N)$ work in
$O(\log N)$ depth by exploiting the associativity of~$\oplus$ to
evaluate sub-ranges independently and compose the results.  On a GPU,
the challenge is implementing this efficiently without excessive
synchronization.

\begin{figure}[!htbp]
\centering
\begin{tikzpicture}[
    >=Stealth,
    scale=0.72,
    blk/.style={draw, rounded corners=2pt, minimum width=0.9cm,
                minimum height=0.5cm, font=\scriptsize, align=center},
    barr/.style={draw, red!70!black, line width=1.5pt, dashed},
    barrlbl/.style={font=\tiny\bfseries, red!70!black},
    arr/.style={->, thick, gray!70!black},
]
  % === TOP PANEL: Hierarchical (PPv3) ===
  \node[font=\small\bfseries] at (6.5, 8.5)
    {(a) Hierarchical scan (PPv3) --- 3 global barriers};

  % Block row
  \foreach \i/\lbl in {0/CTA 0, 1/CTA 1, 2/CTA 2, 3/CTA 3, 4/CTA 4, 5/CTA 5} {
    \node[blk, fill=blue!15] (b\i) at (2.2*\i, 7.5) {\lbl};
  }
  \node[font=\tiny, gray] at (6.5, 7.0) {Step 1: each CTA scans its local chunk};

  % Barrier 1
  \draw[barr] (-0.8, 6.5) -- (12.0, 6.5);
  \node[barrlbl] at (12.5, 6.5) {sync 1};

  % Global scan by CTA 0
  \node[blk, fill=orange!20, minimum width=6cm] (gscan) at (5.5, 5.8)
    {CTA 0 scans all aggregates};
  \foreach \i in {0,...,5} {
    \draw[arr] (b\i.south) -- (gscan.north -| b\i.south);
  }

  % Barrier 2
  \draw[barr] (-0.8, 5.1) -- (12.0, 5.1);
  \node[barrlbl] at (12.5, 5.1) {sync 2};

  % Apply prefix back
  \foreach \i/\lbl in {0/apply, 1/apply, 2/apply, 3/apply, 4/apply, 5/apply} {
    \node[blk, fill=green!12] (a\i) at (2.2*\i, 4.4) {\lbl};
    \draw[arr] (gscan.south -| a\i.north) -- (a\i.north);
  }
  \node[font=\tiny, gray] at (6.5, 3.8)
    {Step 3: each CTA applies global prefix to local results};

  % Barrier 3
  \draw[barr] (-0.8, 3.3) -- (12.0, 3.3);
  \node[barrlbl] at (12.5, 3.3) {sync 3};

  % === BOTTOM PANEL: Decoupled look-back (PPv5) ===
  \node[font=\small\bfseries] at (6.5, 2.2)
    {(b) Decoupled look-back (PPv5) --- no global barriers};

  % Blocks working independently at staggered times
  \foreach \i/\xoff in {0/0, 1/0.6, 2/1.0, 3/0.3, 4/1.4, 5/0.8} {
    \pgfmathsetmacro{\ybase}{0.7}
    \node[blk, fill=blue!15] (d\i) at (2.2*\i, \ybase) {CTA \i};
  }

  % Look-back arrows (only some blocks need them)
  \draw[arr, orange!70!black, bend right=20] (d1.south west) to
    node[below, font=\tiny, orange!70!black] {look-back} (d0.south east);
  \draw[arr, orange!70!black, bend right=20] (d3.south west) to (d2.south east);
  \draw[arr, orange!70!black, bend right=15] (d5.south west) to (d4.south east);

  \node[font=\tiny, gray, align=center] at (6.5, -0.3)
    {Blocks proceed independently; look-back resolves\\inter-block
     dependencies on demand (no global barrier)};
\end{tikzpicture}
\caption[Hierarchical vs.\ decoupled look-back parallel prefix]{%
  Two parallel prefix scan strategies for GPU carry propagation.
  \emph{(a)}~The hierarchical approach (PPv3) requires three global
  \code{grid.sync()} barriers: after local scans, after the global
  aggregate scan (performed by CTA~0 alone), and after applying the
  global prefix back to each CTA's local results.
  \emph{(b)}~The decoupled look-back approach (PPv5) eliminates global
  barriers entirely: each CTA publishes its aggregate as soon as its
  local scan completes and resolves upstream dependencies via a
  lightweight look-back that queries only the nearest unresolved
  predecessor.}
\label{fig:prefix-scan-comparison}
\end{figure}

\paragraph{Contrast with Blelloch-style parallel prefix scans.}

Classical parallel prefix algorithms, most notably the Blelloch
scan \cite{Blelloch1990PrefixSums}, provide a
well-established framework for evaluating associative operators with
$O(\log N)$ depth and $O(N)$ work. Blelloch-style scans proceed in two distinct
phases: an up-sweep (reduce) phase that computes partial aggregates over a
balanced tree, followed by a down-sweep phase that distributes prefix values to
all elements. This structure is well suited to SIMD and shared-memory parallel
models and forms the basis of many CPU and GPU prefix implementations.

However, on GPUs, Blelloch scans impose rigid global synchronization points
between phases. When applied to large limb arrays, this typically requires
either multiple kernel launches or explicit grid-wide synchronization barriers.
As a result, the execution is dominated by synchronization latency rather than
arithmetic throughput, particularly when the scan is embedded inside a larger
iterative algorithm such as high-precision reference orbit computation.

In contrast, the single-pass parallel prefix algorithm
of \cite{merrill2016single} eliminates the global phase separation inherent in
Blelloch scans (\cref{fig:prefix-scan-comparison}).Instead of enforcing a strict up-sweep/down-sweep structure,
each thread block computes its local prefix independently and resolves
inter-block dependencies dynamically using a decoupled look-back mechanism. This
allows blocks to make forward progress asynchronously and avoids global barriers
in the common case.

For high-precision carry propagation, this distinction is critical. Carry chains
often terminate locally, and only a small fraction of limbs depend on distant
predecessors. Blelloch-style scans nevertheless force all blocks to participate
in every global synchronization phase, regardless of whether a true dependency
exists. Merrill--Garland, by contrast, pays synchronization cost only when a
carry must cross block boundaries, yielding substantially better utilization of
GPU resources.

From a numerical perspective, both approaches compute the same associative
prefix over carry transfer operators and are therefore equally correct. The
difference lies entirely in execution strategy. Blelloch scans provide a
deterministic, phase-structured evaluation with predictable synchronization,
while Merrill--Garland offers a latency-tolerant, demand-driven execution model
that is better aligned with the irregular dependency structure of carry
propagation in large arbitrary-precision additions.

In this work, the Merrill--Garland formulation is preferred not because it
changes the arithmetic, but because it minimizes synchronization on parallel
hardware, allowing carry-dominated operations to scale to precisions
where serial or multi-pass approaches become impractical.

\paragraph{Historical note: hierarchical scan implementation.}
An earlier carry propagation implementation
(\code{CarryPropagation\_ABC\_PPv3} in \code{Add\_ABC.h}) uses a
classical hierarchical parallel prefix scan: each CTA computes a local
inclusive scan of its chunk's transfer functions, publishes its
aggregate, then CTA~0 scans the per-CTA aggregates, and finally each
CTA incorporates the global prefix.  This requires three explicit
\code{grid.sync()} barriers within the scan alone (local$\to$publish,
global scan, apply prefix), in addition to the barriers before and
after the scan.  The implementation is retained in the codebase for
reference but is not used in production; the decoupled look-back
approach (PPv5) replaces it with fewer barriers and better scaling.


\subsubsection{Discussion}

By expressing carry propagation as a parallel prefix problem and employing a
single-pass, decoupled look-back strategy, high-precision floating-point
addition becomes compatible with GPU execution. The formulation cleanly
separates numerical correctness from execution strategy and composes naturally
with NTT-based multiplication and fused arithmetic pipelines. As a result,
fully high-precision arithmetic becomes practical for demanding workloads such
as GPU-accelerated reference orbit computation.


\subsection{Number Theoretic Transform and Multiply}
\label{sec:ntt-multiply}

High-precision floating-point multiplication ultimately reduces to multiplying
two large integers (the significands), followed by normalization and exponent
adjustment. A standard way to accelerate large-integer multiplication is to
compute the convolution of digit-limbs using a fast transform. Over the reals,
one would use an FFT; over modular arithmetic, the analogous tool is the
\emph{Number Theoretic Transform} (NTT) \cite{NTT_Wikipedia,NTT_BeginnerGuide2024}.

\subsection{From Large-Integer Multiplication to Convolution}

Let the two nonnegative integers to be multiplied be represented in base $B$ as

\[
A = \sum_{i=0}^{n-1} a_i B^i,\qquad
C = \sum_{i=0}^{n-1} c_i B^i,
\]

with $0 \le a_i, c_i < B$. Their product is

\[
A\cdot C = \sum_{k=0}^{2n-2} \left(\sum_{i=0}^{k} a_i c_{k-i}\right) B^k.
\]

The coefficients

\[
s_k = \sum_{i=0}^{k} a_i c_{k-i}
\]

form the \emph{discrete convolution} of the sequences $(a_i)$ and $(c_i)$. If we
can compute $(s_k)$ quickly, then we can recover the product (with subsequent
carry propagation in base $B$).

\subsection{NTT as an FFT Over a Finite Field}

The NTT computes a discrete Fourier transform, but with all operations performed
modulo a prime $p$ rather than over $\mathbb{C}$. Choose a prime modulus $p$ and
an $N$-th primitive root of unity $\omega \in \mathbb{Z}_p$ such that

\[
\omega^N \equiv 1 \pmod p,\qquad
\omega^k \not\equiv 1 \pmod p \text{ for } 0<k<N.
\]

Then the forward NTT of a length-$N$ sequence $x = (x_0,\dots,x_{N-1})$ is

\[
X_k = \sum_{j=0}^{N-1} x_j \,\omega^{jk} \pmod p,\qquad k=0,\dots,N-1,
\]

and the inverse NTT is

\[
x_j = N^{-1} \sum_{k=0}^{N-1} X_k \,\omega^{-jk} \pmod p,\qquad j=0,\dots,N-1,
\]

where $N^{-1}$ is the multiplicative inverse of $N$ modulo $p$.

Just as with the complex FFT, the transform diagonalises convolution:

\[
\mathrm{NTT}(x \star y) \;=\; \mathrm{NTT}(x)\odot \mathrm{NTT}(y),
\]

where $\odot$ denotes pointwise multiplication and $\star$ denotes cyclic
convolution modulo $x^N-1$. To compute the \emph{linear} convolution needed for
multiplication, we pad both inputs with zeros to length $N \ge 2n$ so that the
cyclic convolution coincides with the linear convolution.

\subsection{Why a Special Prime Helps: \code{MagicPrime}}

A practical NTT requires:

\begin{enumerate}

\item $N$ to be highly composite (typically a power of two) so that
Cooley--Tukey style butterflies apply efficiently;

\item a modulus $p$ such that $N \mid (p-1)$, guaranteeing the existence of an
$N$-th primitive root of unity in $\mathbb{Z}_p$;

\item fast modular multiplication on the target hardware (here, GPUs with
efficient 64-bit integer operations).

\end{enumerate}

The prime used here, denoted \code{MagicPrime}, is chosen specifically to
satisfy these constraints. The critical mathematical implication of using a
prime modulus is that $\mathbb{Z}_p$ is a field, so every nonzero element has a
multiplicative inverse and the NTT is well-defined and invertible when $\omega$
exists.

The most important structural requirement is:

\[
N \mid (p-1).
\]

When $N$ is a power of two, say $N = 2^m$, this becomes

\[
2^m \mid (p-1).
\]

Thus, the larger the power of two dividing $(p-1)$, the larger the supported
transform sizes.

\subsection{The ``Goldilocks'' Form and Power-of-Two Roots}

A particularly convenient choice on 64-bit hardware is a prime of the form

\[
p = 2^{64} - 2^{32} + 1,
\]

often called a ``Goldilocks'' prime. (In the codebase this role is played by
\code{MagicPrime}.) Two consequences make this form attractive:

\paragraph{(1) Large power-of-two factor in $p-1$.}

We have

\[
p - 1 = 2^{64} - 2^{32} = 2^{32}\left(2^{32}-1\right),
\]

so $2^{32}$ divides $(p-1)$. Therefore, for any $N = 2^m$ with $m \le 32$, there
exists an $N$-th root of unity in $\mathbb{Z}_p$. This enables power-of-two NTTs
up to length $2^{32}$ in principle (practically limited by memory and
implementation constraints), which is ample for large-limb convolutions.

\paragraph{(2) Efficient modular arithmetic with 64-bit operations.}

While modular multiplication in $\mathbb{Z}_p$ conceptually involves products up
to $p^2$, this modulus is engineered so reductions can be implemented
efficiently using 64-bit and 128-bit intermediates and/or Montgomery reduction.
The prime is close to $2^{64}$, which aligns well with native unsigned integer
ranges, and its special structure admits reduction strategies that avoid slow
division.

\subsection{NTT-Based Multiplication at a High Level}

Given limb sequences $(a_i)$ and $(c_i)$:

\begin{enumerate}

  \item Choose $N$ as a power of two with $N \ge 2n$, and choose a primitive
  $N$-th root $\omega \in \mathbb{Z}_p$.

  \item Zero-pad both sequences to length $N$ (interpreting limbs as elements of
  $\mathbb{Z}_p$).

  \item Compute $A_k = \mathrm{NTT}(a)_k$ and $C_k = \mathrm{NTT}(c)_k$.

  \item Compute pointwise products $P_k = A_k \cdot C_k \pmod p$.

  \item Compute $p_j = \mathrm{NTT}^{-1}(P)_j$ to obtain the convolution
  coefficients modulo $p$.

\end{enumerate}

Finally, because the convolution coefficients are computed modulo $p$, we must
ensure they represent the \emph{true integer} convolution values, not values
wrapped modulo $p$. This is achieved by choosing the limb base $B$ and transform
length $N$ so that each exact coefficient $s_k$ satisfies

\[
0 \le s_k < p,
\]

(or more generally, can be reconstructed from one or more moduli). With a single
sufficiently large prime (as is typical with a 64-bit Goldilocks prime and
appropriately sized limbs), the coefficients can be recovered directly and then
normalized via carry propagation in base $B$.

\medskip

This section establishes the mathematical foundation: \code{MagicPrime} is
selected so that large power-of-two NTTs exist (because $2^m \mid p-1$) and so
that modular arithmetic is efficient on 64-bit GPU hardware. Subsequent sections
will describe how this is specialized for high-precision floating-point
significands, including limb packing, Montgomery-domain multiplication, twiddle
scheduling, and normalization/carry handling.

\subsection{NTT in \FractalShark{}}

The NTT implementation in \FractalShark{} does not prioritize
memory layout or optimized coalescing, which likely leaves performance on the
table.  This section describes how the NTT-based
multiplication is realized in practice, following the structure of
\FractalShark{}'s CUDA implementation. The overall goal is to compute the exact
convolution of two large significand arrays using modular arithmetic modulo
\code{MagicPrime}, and then to recover a canonical high-precision result via
normalization and carry propagation.

\subsubsection{Plan Construction and Parameter Selection}

Before launching any GPU kernels, an NTT ``plan'' is constructed. The plan
determines:

\begin{itemize}

  \item the coefficient bit-width $b$ used to pack the original 32-bit limbs
  into NTT coefficients,

  \item the packed coefficient length $L$,

  \item the transform size $N$, chosen as a power of two with $N \ge 2L$.

\end{itemize}

Correctness requires that no convolution coefficient overflow the prime modulus.
If the coefficients are bounded by $2^b$ and the transform length is $N$, then a
conservative bound on the largest convolution term is

\[
\max_k s_k \;\le\; N \cdot 2^{2b}.
\]

The plan builder enforces

\[
2b + \log_2 N + \delta \;\le\; 64,
\]

for a small safety margin $\delta$, ensuring that all exact convolution values
lie strictly below \code{MagicPrime}. This guarantees that the modular
convolution coincides with the true integer convolution.

\subsubsection{Roots of Unity and Montgomery Domain}

The CUDA setup phase precomputes all roots of unity required for the radix--2
NTT. For each stage $s$, a primitive $2^s$-th root $\omega_s$ and its inverse
are generated, along with:

\begin{itemize}

  \item powers of a twist root $\psi^i$ and $\psi^{-i}$,

  \item the modular inverse $N^{-1}$.

\end{itemize}

All constants are stored in Montgomery representation \cite{Montgomery1985Modular}. As a result, every
modular multiplication inside the NTT butterflies is implemented as a Montgomery
multiply, avoiding explicit modular reduction by division.

\subsubsection{Packing, Twisting, and Forward NTT}

The input significands are initially stored as arrays of 32-bit limbs. These are
packed into base-$2^b$ coefficients

\[
A(x) = \sum_{i=0}^{L-1} a_i x^i,\qquad
C(x) = \sum_{i=0}^{L-1} c_i x^i,
\]

with $0 \le a_i, c_i < 2^b$. During packing, each coefficient is:

\begin{enumerate}

  \item mapped into $\mathbb{Z}_p$,

  \item multiplied by the twist factor $\psi^i$,

  \item converted into Montgomery form.

\end{enumerate}

After packing, an in-place radix--2 forward NTT is performed using the
decimation-in-time (DIT) formulation.  DIT requires the input to be in
\emph{bit-reversed} order: each element at index~$i$ must appear at
index~$\mathrm{bitrev}(i)$, where $\mathrm{bitrev}$ reverses the
$\log_2 N$ least-significant bits.  This ensures that the iterative
butterfly stages can read and write elements at naturally strided
offsets, enabling coalesced memory access patterns on the GPU.

\FractalShark{} fuses the forward bit-reversal into the packing step:
rather than writing packed/twisted coefficients to position~$i$ and then
performing a separate permutation pass, the packing kernel computes
$j = \mathrm{bitrev}(i)$ via the hardware \code{\_\_brev} intrinsic
and scatter-writes each coefficient directly to position~$j$.
This eliminates one full random-access permutation pass over the
$N$-element array, saving both memory traffic and a
\code{grid.sync()} barrier.

For the inverse NTT, a separate in-place bit-reversal is still
performed via \code{BitReverseInplace64\_GridStride}, which
permutes all active arrays simultaneously (\code{SevenWay} when NR is
enabled, \code{ThreeWay} otherwise).  Each thread swaps
elements $i$ and $\mathrm{bitrev}(i)$ (with the $i < j$ guard ensuring
each pair is swapped exactly once).

\paragraph{DIT vs.\ DIF and the bit-reversal tradeoff.}
The DIT formulation requires bit-reversal on the \emph{input}; the
alternative decimation-in-frequency (DIF) formulation produces output in
bit-reversed order, requiring the permutation \emph{after} the
transform.  A mixed strategy---DIT for the forward transform and DIF for
the inverse (or vice versa)---can cancel the remaining inverse
permutation entirely, since the bit-reversed output of the forward
transform would already be in the correct input order for a DIF inverse.
This ``self-sorting'' approach would eliminate the last standalone
permutation pass.

\FractalShark{} currently uses DIT for both forward and inverse
transforms; the forward bit-reversal has been fused into packing, but
the inverse still incurs a separate permutation.  Adopting a DIT/DIF
mixed strategy would remove this remaining pass and its associated
memory traffic, at the cost of a different butterfly indexing pattern in
the inverse transform.

After bit-reversal, iterative butterfly stages are applied:

\[
(u, v) \;\mapsto\; (u + \omega v,\; u - \omega v),
\]

with all operations performed modulo \code{MagicPrime}. A grid-stride loop is
used so that threads collectively cover all $N$ coefficients, independent of the
exact grid dimensions.

\subsubsection{Multiway Pointwise Multiplication}

In the transform domain, convolution reduces to pointwise multiplication. To
amortize the cost of the NTT, the implementation computes three related
convolutions simultaneously using a standard three-multiply decomposition.
Conceptually, if

\[
X = X_r + iX_i,\qquad Y = Y_r + iY_i,
\]

then the products

\[
X_rY_r,\quad X_iY_i,\quad (X_r+X_i)(Y_r+Y_i)
\]

are sufficient to reconstruct both real and imaginary components. Accordingly,
three frequency-domain products are computed per transform index:

\[
\widehat{XX}_k,\quad \widehat{YY}_k,\quad \widehat{XY}_k,
\]

each via a Montgomery modular multiplication modulo \code{MagicPrime}.

\subsubsection{Inverse NTT and Untwisting}

Following pointwise multiplication, inverse radix--2 NTTs are applied using the
inverse roots of unity. The inverse transform yields coefficients still in
Montgomery form and still containing the twist factor. These are corrected by:

\begin{enumerate}
\item multiplying by $\psi^{-i}$,
\item multiplying by $N^{-1}$,
\item converting out of Montgomery representation via
      $\operatorname{FromMontgomery}(v) = \operatorname{MontMul}(v,\,1)$.
\end{enumerate}

\noindent
In practice these three corrections are fused into a single grid-stride pass
(\code{Untwist\-Scale\-From\-Mont\_\allowbreak 3Way\_\allowbreak Grid\-Stride}):
for each NTT index~$i$ and each product channel, two successive Montgomery
multiplications apply the precomputed inverse twist factor
$\psi^{-i}\bmod p$ and the precomputed inverse transform length
$N^{-1}\bmod p$, after which a final \code{FromMontgomery} call reduces
the result out of the Montgomery domain.  Because all three operands
($\psi^{-i}$, $N^{-1}$, and the constant~$1$ used by
\code{FromMontgomery}) are known at kernel-launch time, the root-table
setup precomputes them and stores them in device memory alongside the
forward roots.

\subsubsection{Unpacking and Normalization}
\label{subsec:unpack-normalize}

The untwisted, scaled coefficients are elements of $\mathbb{Z}/p\mathbb{Z}$,
each fitting in 64~bits.  Converting them back to a positional
base-$2^{32}$ limb representation requires distributing each
coefficient's contribution across adjacent digit boundaries---a step
called \emph{unpacking}.  The kernel \code{UnpackPrimeToFinal128}
iterates over target digit positions: for each digit~$j$, it collects the
overlapping bit ranges from the NTT coefficients that straddle that
boundary, testing whether each coefficient represents a positive or
negative residue (values above $(p-1)/2$ are treated as negative),
and accumulating the results into 128-bit (low, high) pairs.

Once unpacking completes, the 128-bit pairs may contain inter-digit
carries that span multiple positions.  A multi-phase carry-propagation
pass resolves them.  Two rounds of digit-by-digit ripple carry first
reduce the wide accumulators into 32-bit digits with at most single-bit
residual carries.  If any residual carries survive both rounds, a
decoupled-look-back parallel prefix scan
(\cref{sec:hp-add}) resolves them in a single
additional pass.  The result is a fully normalized base-$2^{32}$ limb
array that is written back into the destination \code{HpSharkFloat}
buffer together with the product's sign and exponent (computed from the
input operands' exponents).

\subsection{Pipeline synthesis}
\label{subsec:ntt-pipeline-synthesis}

Taken together, the stages described in this section form a single
end-to-end GPU kernel invocation for high-precision multiplication.
The forward path packs limbs into the NTT domain, applies twist
factors, and performs a radix-2 forward NTT; pointwise Montgomery
multiplications then produce the convolution in the frequency domain.
The inverse path applies an inverse NTT, removes the twist factors,
scales by $N^{-1}$, and exits the Montgomery domain.  Finally,
unpacking and carry normalization convert the prime-field
coefficients back into a positional limb representation.

Because multiple independent products (at minimum the three real-times-real
channels $XX$, $YY$, $XY$ needed for complex multiplication, and up to
seven channels when Newton--Raphson derivatives are included) share the
same NTT length and modular-arithmetic parameters, every stage is
fused into a single grid-stride loop and all channels advance through
each barrier together.  This batching strategy, described in
\cref{subsec:amortize-sync}, amortises the cost of grid-wide
synchronisation across all simultaneously computed products and is
central to the practical performance of GPU-resident reference-orbit
computation.

\subsection{Amortizing Synchronization via Simultaneous Products}
\label{subsec:amortize-sync}

In the presented GPU implementation, a dominant cost of large-scale NTT-based
multiplication is not the modular arithmetic itself, but the required grid-wide
synchronization. Each major phase of the pipeline— packing and twisting of
inputs, forward NTT, pointwise multiplication, inverse NTT, and untwisting with
normalization—requires that all threads observe a consistent global state. This
is enforced using explicit grid-level barriers (e.g.,
\code{cooperative\_groups::grid\_group::sync()}), whose latency grows with
grid size and is largely independent of the amount of arithmetic performed
between barriers.

To reduce the amortized cost of these synchronization points, the implementation
computes three related convolution products, $X\cdot Y$, $X^2$, and $Y^2$,
within a single NTT pipeline. All three products share the same forward
transforms, inverse transforms, twiddle scheduling, temporary buffers, and
synchronization boundaries. As a result, the fixed cost of grid-wide
coordination and memory visibility is incurred once, while producing three
mathematically independent convolution results. The additional arithmetic
required for the extra pointwise products consists only of a small number of
Montgomery multiplications per frequency index and is negligible compared to the
cost of global synchronization and memory traffic.

At the code level, this strategy is realized by a fused front-end that packs the
input limb arrays, applies twist factors, and converts values into Montgomery
form while producing multiple frequency-domain streams in a single pass. A
single grid-stride pointwise multiplication phase then computes the three
products concurrently. These are followed by a single multiway inverse radix--2
NTT, after which a unified untwisting, scaling by $N^{-1}$, and conversion out
of Montgomery representation are performed. Explicit grid-wide barriers appear
only at true phase boundaries, serving as producer--consumer separators between
global-memory stages rather than as fine-grained synchronization.

This design is well matched to the target workload. In high-precision reference
orbit computation, both cross terms and squared terms arise naturally and
repeatedly. By folding these operations into a single multiway NTT, the
implementation reduces the total number of kernel phases and synchronization
events, increases arithmetic intensity per barrier, and improves overall GPU
utilization. Consequently, the effective performance of the multiplication
pipeline is governed primarily by arithmetic throughput and memory bandwidth,
rather than by synchronization overhead.  A concrete application of this
principle appears in \cref{subsec:ff-inner-gpu}, where Newton--Raphson
derivative products are interleaved into the orbit-multiplication pipeline
with zero additional synchronization cost.


\subsection{Checksum-Guided Debugging via Host--Device Cross-Validation}
\label{sec:checksums}

Correctness bugs in GPU high-precision arithmetic are difficult to
localize: a single off-by-one carry, lane-misaligned shuffle, or missing
synchronization can corrupt results far downstream, often without an obvious
local symptom. To shorten the feedback loop, we instrument the arithmetic
pipeline (\cref{sec:hp-add,sec:ntt-multiply}) with lightweight \emph{stage checksums} that can be computed on the GPU,
computed independently on the host reference implementation, and compared to
pinpoint the first divergence. This converts a diffuse ``wrong final answer''
into a small set of candidate stages, reducing the search space
during development.  These checksums are disabled in production builds to avoid
any performance impact but can be enabled selectively during debugging.

\subsubsection{Stage Checksums as Homomorphic Summaries}

Let $A \in (\mathbb{Z}/2^w\mathbb{Z})^N$ denote a limb array representing an
intermediate mantissa or scratch buffer (e.g., aligned addends, raw limbwise
sums, carry descriptors, normalized limbs). We associate to each such array a
64-bit checksum
\begin{equation}
  H(A) \in \{0,1\}^{64},
\end{equation}
computed deterministically over the ordered limb sequence. In our implementation
we employ 64-bit rolling checksums (e.g., CRC-64 or Fletcher-64) because they are
(1) inexpensive relative to high-precision arithmetic, (2) order-sensitive, and
(3) easily reducible in parallel.

Each \emph{debug checkpoint} records a tuple
\begin{equation}
  S_k = \bigl(\code{Purpose}_k,\, \code{Depth}_k,\, \code{CallIndex}_k,\,
  \code{Conv}_k,\, H(A_k)\bigr),
\end{equation}
where $\code{Purpose}$ identifies which logical stage or buffer is being
summarized (e.g., \code{ADigits}, \code{BDigits}, \code{FinalAdd1},
\code{Result\_offsetXX}, etc.), and the remaining metadata disambiguates
dynamic execution contexts such as recursion depth or multiple invocations of
the same primitive within a larger operation. This metadata ensures that host
and device checkpoints can be matched without ambiguity even when the same
kernel path is exercised repeatedly in a pipeline.

\subsubsection{Parallel Computation with Order Preservation}

The checksum must match between host and device, so the GPU computation must be
\emph{equivalent} to the sequential definition. This requires two properties:
\begin{enumerate}
  \item \textbf{Deterministic segmentation}: the array is partitioned into
  contiguous chunks assigned to threads in a deterministic manner.
  \item \textbf{Order-preserving reduction}: partial checksums are combined in
  the same left-to-right order as the sequential traversal.
\end{enumerate}

To make this concrete, consider a checksum family whose internal state can be
represented as a small tuple (e.g., for Fletcher-64, $(s_1,s_2)$ with implicit
length). Each thread processes a contiguous chunk $A[\ell:r)$ and produces a
local state $\sigma_j$. To recover the global checksum, we require an
associative \emph{combine} operator $\oplus$ such that
\begin{equation}
  \sigma(A[0:n)) \;=\; \sigma(A[0:k)) \;\oplus\; \sigma(A[k:n)).
\end{equation}
Associativity enables hierarchical reduction, while order preservation requires
that the reduction tree respects the left-to-right concatenation order. On the
GPU, we implement this reduction hierarchically:
\begin{enumerate}
  \item Each thread computes $\sigma_j$ for its chunk.
  \item Within each warp, we compute an order-preserving prefix/reduction using
  warp shuffles and $\oplus$.
  \item Warp results are combined in increasing warp order to form a per-block
  checksum.
  \item Block results are reduced across the grid using repeated pairwise
  left-to-right combinations, yielding a single 64-bit checksum.
\end{enumerate}
This structure ensures that the device result is \emph{definitionally identical}
to the host's sequential traversal, modulo the checksum's arithmetic rules.

\subsubsection{Host Reference and Cross-Comparison}

On the host, we compute the same checksum $H(A_k)$ over the corresponding
reference buffers produced by the MPIR/GMP-based implementation (or a trusted
CPU implementation). Because both sides share the identical stage labeling
(\code{Purpose}, \code{Depth}, \code{CallIndex}, \code{Conv}), we can
compare checkpoint sequences as keyed records:
\begin{equation}
  \Delta_k \;=\;
  \bigl(H_{\mathrm{gpu}}(A_k) \stackrel{?}{=} H_{\mathrm{host}}(A_k)\bigr).
\end{equation}

A mismatch at checkpoint $k$ implies that the first divergence occurred at or
before the computation that produces $A_k$. By inserting checkpoints at
strategically chosen boundaries (e.g., after exponent alignment, after signed
limbwise accumulation, after carry resolution, after normalization), we obtain a
coarse-to-fine localization mechanism. In practice this behaves like a binary
search over the pipeline: we start with a small number of high-level checkpoints
and then refine around the earliest failing stage by adding more granular
summaries (e.g., per-buffer or per-substep).

\subsubsection{Practical Debugging Workflow}

The checksum mechanism supports a fast iterative workflow:
\begin{enumerate}
  \item \textbf{Instrument}: enable checksum capture for selected stages and
  record the resulting $(\code{key}, H)$ tuples from GPU execution.
  \item \textbf{Replay on host}: run the corresponding host reference path over
  the same inputs, generating the same keyed tuples.
  \item \textbf{Diff}: locate the earliest key for which
  $H_{\mathrm{gpu}} \ne H_{\mathrm{host}}$.
  \item \textbf{Refine}: insert additional checkpoints in the immediate
  neighborhood of the first mismatch (e.g., split ``carry propagation'' into
  descriptor publish, look-back resolution, digit transfer, and final writeback).
\end{enumerate}

This approach is especially effective for bugs that only manifest at scale:
race conditions, mis-specified synchronization, and rare carry-chain corner
cases. Even when the final numerical error is small or intermittent, checksum
mismatches provide a reproducible signature of divergence.

\subsubsection{Why Checksums Help for High-Precision Arithmetic}

High-precision arithmetic pipelines are composed of large, structured arrays
whose semantics evolve across stages (aligned limbs $\rightarrow$ raw sums
$\rightarrow$ carry-resolved digits $\rightarrow$ normalized mantissa). A
checksum provides a compact witness of correctness for each stage without
requiring full array dumps or expensive per-element comparisons. While a checksum
does not identify the exact index of the first wrong limb, it reliably answers a
more valuable question during early debugging: \emph{which transformation first
produced an incorrect state?} Once localized, targeted assertions, partial dumps,
or reduced test cases can isolate the exact defect with far less effort.

Taken together, stage checksums act as a low-overhead correctness scaffold that
bridges device execution and a trusted host reference, enabling rapid fault
localization in complex GPU high-precision arithmetic kernels.

\subsection{Usage in \FractalShark{}}
At the application layer, the GPU reference orbit backend is exposed through
the persistent \emph{combo} object described earlier and a simple three-stage
lifecycle:

\begin{enumerate}
  \item \code{InitHpSharkReferenceKernel}: allocates persistent device/host
        state and binds the orbit parameters and launch configuration.
  \item \code{InvokeHpSharkReferenceKernel}: advances the orbit in bounded
        batches of at most \(\code{MaxOutputIters}\), streaming results through
        \code{OutputIters}.
  \item \code{ShutdownHpSharkReferenceKernel}: frees the persistent resources.
\end{enumerate}

Because the precision must be fixed at compile time for the \code{
SharkFloatParams} specialization, \code{DispatchByPrecision} rounds the
requested precision to a power of two and chooses from a fixed supported set
(\(\{256,512,\dots,524288\}\) bits). Each invocation appends the emitted
\code{OutputIters} records into the CPU-side \code{PerturbationResults}, while
periodicity and escape are reported through \code{PeriodicityStatus}. From the
rest of \FractalShark{}'s perspective, the backend therefore appears as a
batch-driven, fixed-precision reference-orbit engine whose internal arithmetic
machinery remains encapsulated behind this host-facing lifecycle.

The same GPU high-precision arithmetic powers another major subsystem.  The
next chapter describes the Feature Finder, which applies Newton--Raphson
iteration at arbitrary precision to automatically locate periodic points
embedded deep within the Mandelbrot set.
